


\lstdefinelanguage{xquery}{
  keywords={for, in, collection, doc, return, to, text, order, by, count, descending, let, where, avg, if, then, else, satisfies, contains, and, every, variable, order, while, do, else, case, break, true, false, document, element},
  keywordstyle=\color{BlueViolet}\ttfamily,
%  ndkeywords={class, export, boolean, throw, implements, import, this},
%  ndkeywordstyle=\color{darkgray}\bfseries,
  identifierstyle=\color{black},
%  sensitive=false,
  morecomment=[s]{(:}{:)},
  commentstyle=\color{purple}\ttfamily,
  stringstyle=\color{DarkOliveGreen}\ttfamily,
  morestring=[b]"
}

\section{Préliminaires}

\lstset{language=xquery,basicstyle=\normalsize}

\begin{frame}[allowframebreaks,t]{Modèle de données et expressions}

XQuery travaille sur des \alert{sequences}, notion très générale.

\vfill

Pour ce cours: on considère un cas spécial de séquence, la \alert{collection}, constituée de documents XML.

\vfill

Une \alert{expression} XQuery prend en entrée une collection et renvoie une collection.

\vfill

Deux fonctions permettent d'accéder aux collections: 

\begin{itemize}
  \item \lstinline!doc()! prend en entrée l'URI d'un document XML et renvoie une collection-singleton 
contenant \alert{le nœud racine} du document.

  \item \lstinline!collection()! prend en entrée l'URI d'une collection et 
renvoie une séquence composée \alert{des nœuds racine} des documents de la collection.
\end{itemize}

Dans ce qui suit: exemples sur la collection \texttt{movies} avec l'ensemble des films de Webscope.

\end{frame}


\begin{frame}[fragile]
\frametitle{XPath et au-delà}

Toute expression XPath est une requête XQuery. Voici une expression XPath/XQuery
recherchant tous les films parus en 2005.

\begin{lstlisting}[language=myXQuery]
collection('movies')/movie[year=2005]/title
\end{lstlisting}

Le résultat est un \alert{séquence} de nœuds \texttt{title}\,::
\begin{lstlisting}[language=myXML]
<title>A History of Violence</title>
<title>Match Point</title>
\end{lstlisting}

\begin{remark}
L'expression XPath est évaluée pour chaque \domtype{Document} de la séquence  \texttt{collection('movies')}.
\end{remark}

\end{frame}


\begin{frame}[fragile]
\frametitle{Constructeurs}

XQuery permet la construction de nouveaux éléments, dont le contenu
peut mêler des éléments littéraux, des valeurs, ou le résultat d'autres requêtes
(composition!).

\begin{lstlisting}[language=myXQuery]
<titles>
  {collection('movies')//title}
</titles>
\end{lstlisting}
 

\begin{remark}
Une expression $e$ doit être encadréee par des accolades
\{\} pour être reconnue. 
\end{remark}
\end{frame}

\begin{frame}[containsverbatim]{Variables}

Des variables peuvent être utilisées dans les expressions XQuery. Par exemple;

\begin{lstlisting}
<employee empid="{$id}">
   <name>{$name}</name>
   {$job} 
   <deptno>{$deptno}</deptno> 
   <salary>{$SGMLspecialist+100000}</salary>
</employee>
\end{lstlisting}

où les variables \lstinline!$id!, \lstinline!$name!, \lstinline!$job!, \lstinline!$deptno! and \lstinline!$SGMLspecialist! sont instanciées (une valeur est associée à chaque variable). \prof{par exemple avec un let ou au sein d'un for}

\end{frame}


\section{L'essentiel\;: la clause FLOWR}

\begin{frame}[containsverbatim,allowframebreaks]{Expressions FLWOR \prof{prononcer ``flower''}}

\begin{itemize}
\item  4 clauses : 
\begin{itemize}
\item  For (itération sur une séquence), 
\item  Let (définition ou instanciation d'une variable), 
\item  Where (prédicat), 
\item  Order by (tri du résultat), and 
\item  Return (contructeur de résultat).
\end{itemize}
\item  Dans l'esprit du SELECT-FROM-WHERE de SQL
\end{itemize}



\framebreak

Dans l'une de ses formes les plus simples, FLW0R est une alternative à XPath.

\begin{lstlisting}
//actor[birth_date>=1960]/last_name
\end{lstlisting}

est équivalent à l'expression (requête) XQuery suivante :

\begin{lstlisting}
let $year:=1960 
for $a in doc("SpiderMan.xml")//actor 
where $a/birth_date >= $year 
return $a/last_name
\end{lstlisting}

\prof{Evidemment, toutes les expressions XQuery ne peuvent pas être écrites en XPath}

\framebreak

Les clauses \lstinline!for! et \lstinline!let!.

\bigskip

Ces deux clauses instancient des variables.

\bigskip

\lstinline!for! instancie une variable en lui faisant prendre (successivement) les valeurs des items d'une séquence en entrée\\
\medskip
 Par exemple : \lstinline!for $x in /company/employee! \\
instancie la variable \lstinline!x! en lui faisant successivement prendre pour valeur \alert{chacun} 
des éléments \lstinline!employee! de la séquence issue de l'évaluation de \lstinline!/company/employee!.

\bigskip

\lstinline!let! instancie une variable en lui faisant prendre la valeur de la séquence complète\\
\medskip
 Par exemple : \lstinline!let $x := /company/employee!\\
instancie la variable \lstinline!x! en lui faisant prendre la valeur de \alert{la séquence} issue 
de l'évaluation de \lstinline!/company/employee!.

\framebreak

Instanciation d'une variable :
\begin{lstlisting}
let $years:=(1960,1978)
return $years
\end{lstlisting}


Parcours d'une séquence :
\begin{lstlisting}
for $i in (1,2,3)
return 
<chantons>
  { $i } kilometre(s) a pieds,
  Ca use, ca use
  { $i } kilometre(s) a pieds,
  Ca use les souliers
</chantons>
\end{lstlisting}
 
\framebreak

La séquence peut être obtenue par évaluation d'une expression XQuery (ci-dessous XPath)

\bigskip

Instanciation d'une variable :
\begin{lstlisting}
let $speechs := doc("hamlet.xml")//SPEECH
return $speechs
\end{lstlisting}

\bigskip

Parcours d'une séquence :
\begin{lstlisting}
for $speech in  doc("hamlet.xml")//SPEECH
return $speech
\end{lstlisting}

\prof{meme résultat mais comportement different. Expliquer.}

\bigskip
Pour travailler le résultat :

\begin{lstlisting}
for $speech in  doc("hamlet.xml")//SPEECH
return $speech/SPEAKER
\end{lstlisting}

\framebreak

Un autre example :

\begin{lstlisting}
let $born := ('born, born, born')
let $alive := (' to be alive')
for $i in (1,2,3,4,5,6)
   return 
	<refrain numero="{$i}">
	You see you were
	{$born} {$alive}
	You see you were
	{$born}
	born {$alive}
	</refrain>
\end{lstlisting}

\begin{exoblock}
Que renvoie cette expression XQuery ?
\end{exoblock}
\prof{monter correction sous eXist}
\framebreak

Un autre example :

\begin{lstlisting}
for $i in (-3,-2,-1,0,1,2,3)
let $j := $i + 1
return <suivant>{$j} est le nombre suivant {$i}</suivant>
\end{lstlisting}

\begin{exoblock}
Que renvoie cette expression XQuery ?
\end{exoblock}

\prof{monter correction sous eXist}

\framebreak

La clause \lstinline!where! est assez similaire à son synonyme SQL. Elle permet de filtrer un item en fonction d'une condition.

\bigskip

\begin{lstlisting}
for $line in doc("hamlet.xml")//SPEECH/LINE
where contains($line,"To be, or not to be")
return $line/parent::SPEECH/parent::SCENE/TITLE
\end{lstlisting}


\begin{lstlisting}[language=XML]
<TITLE>SCENE I.  A room in the castle.</TITLE>
\end{lstlisting}

\framebreak 

\begin{lstlisting}
for $line in doc("hamlet.xml")//SPEECH/LINE
where contains($line,"heart")
return $line/parent::SPEECH/child::SPEAKER
\end{lstlisting}

\medskip
 
Pour chaque élément (stocké dans \lstinline!$line!) de \lstinline!doc("hamlet.xml")//SPEECH/LINE! \\
satisfaisant \lstinline!contains($line,"heart")! \\
retourner \lstinline!$line/parent::SPEECH/child::SPEAKER!

\framebreak

\prof{testé sous eXist}
\begin{lstlisting}
for $line in doc("hamlet.xml")//SPEECH/LINE
where contains($line,"heart")
return $line/parent::SPEECH/child::SPEAKER
\end{lstlisting}

\begin{lstlisting}[language=XML,basicstyle=\tiny]
<SPEAKER>FRANCISCO</SPEAKER>
<SPEAKER>KING CLAUDIUS</SPEAKER>
<SPEAKER>KING CLAUDIUS</SPEAKER>
<SPEAKER>KING CLAUDIUS</SPEAKER>
<SPEAKER>KING CLAUDIUS</SPEAKER>
<SPEAKER>KING CLAUDIUS</SPEAKER>
<SPEAKER>KING CLAUDIUS</SPEAKER>
<SPEAKER>HAMLET</SPEAKER>
<SPEAKER>LAERTES</SPEAKER>
<SPEAKER>OPHELIA</SPEAKER>
<SPEAKER>HAMLET</SPEAKER>
<SPEAKER>HAMLET</SPEAKER>
<SPEAKER>HAMLET</SPEAKER>
<SPEAKER>HAMLET</SPEAKER>
<SPEAKER>LORD POLONIUS</SPEAKER>
<SPEAKER>HAMLET</SPEAKER>
<SPEAKER>KING CLAUDIUS</SPEAKER>
...
\end{lstlisting}

\framebreak


\struct{where} est comme en SQL, mais... attention, documents semi-structurés. 

\vfill

``Tous les films dirigés par W. Allen''

\begin{lstlisting}[language=myXQuery]
for $m in collection("movies")/movie
where $m/director/last_name="Allen"
return $m/title
\end{lstlisting}

Ca ressemble à une requête SQL. Attention aux différences suivantes\,:
\begin{enumerate}
  \item si un chemin n'xiste pas, pas d'erreur de typage\,: le résultat est \struct{false}.

  \item si un chemin retourne plusieurs nœuds, le résultat est vrai 
  si \alert{au moins} une des comparaisons retourne \texttt{true}.

\end{enumerate}

\vfill

``Tous les films avec Kirsten Dunst'' (note: plusieurs acteurs dans un films!)
\begin{lstlisting}[language=myXQuery]
for $m in collection("movies")/movie
where $m/actor/last_name="Dunst"
return $m/title
\end{lstlisting}

\end{frame}


\begin{frame}[fragile]
\frametitle{La clause \struct{return}}

\struct{return} est obligatoire et sert à construire un ``item'' résultat. 
Un tel item est construit pour chaque lien (\textit{binding}) de la
variables par la clause \struct{for}.


\begin{lstlisting}[language=myXQuery]
for $m in collection("movies")/movie
let $d := $m/director
where $m/actor/last_name="Dunst"
return 
  <div>
   {$m/title/text(), "directed by",  
        $d/first_name/text(),  $d/last_name/text()}, 
     with
     <ol>
      {for $a in $m/actor
        return <li>{$a/first_name, $a/last_name, 
                         " as ",  $a/role}</li>
      } 
     </ol>
  </div>
\end{lstlisting}

\end{frame}
\begin{frame}[containsverbatim,allowframebreaks]{Ordonner le résultat}
La clause \lstinline!order by! permet de spécifier l'ordre de tri du résultat.

\begin{lstlisting}
for $p in doc("hamlet.xml")//PERSONA
order by $p
return $p
\end{lstlisting}


\begin{lstlisting}[language=XML,basicstyle=\tiny]
<PERSONA>A Captain.</PERSONA>
<PERSONA>A Gentleman</PERSONA>
<PERSONA>A Priest. </PERSONA>
<PERSONA>BERNARDO</PERSONA>
<PERSONA>CLAUDIUS, king of Denmark. </PERSONA>
<PERSONA>CORNELIUS</PERSONA>
<PERSONA>English Ambassadors. </PERSONA>
<PERSONA>FORTINBRAS, prince of Norway. </PERSONA>
<PERSONA>FRANCISCO, a soldier.</PERSONA>
<PERSONA>GERTRUDE, queen of Denmark, and mother to Hamlet. </PERSONA>
<PERSONA>GUILDENSTERN</PERSONA>
<PERSONA>Ghost of Hamlet s Father. </PERSONA>
<PERSONA>HAMLET, son to the late, and nephew to the present king.</PERSONA>
<PERSONA>HORATIO, friend to Hamlet.</PERSONA>
<PERSONA>LAERTES, son to Polonius.</PERSONA>
...
<PERSONA>VOLTIMAND</PERSONA>
\end{lstlisting}


\end{frame}

